\def\Ponderation{40}
\def\SigleNumero{STT-1000}
\def\NomCours{Probabilités et statstique}
\def\DateEvaluation{}
\def\HeureEvaluation{}
\def\Session{}
\def\NomEvaluation{Examen 1 - Reprise}

% Ne rien mettre entre les accolades si l'enseignant (ou l'enseignante) est aussi le coordonnateur (ou la coordonnatrice)
\def\Coordonnatrice{}
\def\Coordonnateur{}

% Ne rien mettre entre les accolades s'il y a plus d'une section
\def\Enseignant{}
\def\Enseignante{}

% Ne rien mettre entre les accolades s'il s'agit d'un cours à section unique
% Mettre le nom de chacun des enseignants et enseignantes entre les accolades
% Une case à cocher sera générée pour chacune des sections
\def\SectionA{Anne-Sophie Charest}
\def\SectionB{Julien Miron}
\def\SectionC{}
\def\SectionD{}

\newcommand{\affichepoints}{}     % ← AVEC points
% \renewcommand{\affichepoints}{on} % ← SANS points (non vide)


\def\Directives{
\begin{enumerate}
\item Matériel autorisé:
\begin{itemize}
\item Calculatrice scientifique {\bf autorisée par la faculté des sciences et de génie};
\item Le feuillet de formules fourni avec l'examen. Celui-ci inclut la table de la loi normale. Veuillez ne rien écrire sur le feuillet de formules.
\end{itemize}
\item Assurez-vous que cette évaluation comporte \numquestions ~exercices répartis sur \numpages{}~pages, pour un total de \numpoints{}~points.\
\item \textbf{Veuillez ne pas dégrafer votre examen.}
\item Vous devez\textbf{ justifier toutes vos réponses} par une démarche claire, sauf pour les questions à choix de réponse.\
\item Pour les questions à choix de réponse et les vrais ou faux, vous pouvez mettre un X, un crochet ($\checkmark$) ou remplir la case.\
\item Vous devez laisser tout ce qui ne servira pas durant l'examen à l'avant de la salle AVANT de rejoindre votre place.\
\item Vous n'avez droit à AUCUN appareil électronique en votre possession durant l'examen.
\item Une fois assis à votre place, vous devez demandez l'autorisation pour vous lever, même si l'examen n'est pas encore commencé, sauf lorsque vous avez terminé l'examen.
\item Dans les 5 dernières minutes de l'examen, il sera interdit de se lever avant que TOUTES les copies n'aient été ramassées.
\end{enumerate}
\begin{itemize}
\item[$\rightarrow$] Le non-respect des consignes 7-8-9 entraînera une pénalité de 10\%.
\end{itemize}
}



\documentclass{Evaluation}
\usepackage{multicol,graphicx}

\printanswers

\begin{document}
\begin{questions}

%%%%%%%%%%%%%% MODULE 1 - PARTIE 1 %%%%%%%%%%%%%%%%%%
\question[4] 

On sait que \(P(A)=0{,}3\), \(P(B)=0{,}2\) et \(P(A^c \cap B)=0{,}1\).
Calculez \(P(A \cap B)\).

\begin{solution}
On décompose l’événement \(B\) :
\[
B = (A \cap B) \cup (A^c \cap B),
\]
où les deux événements sont disjoints. Par additivité des probabilités,
\[
P(B) = P(A \cap B) + P(A^c \cap B).
\]

On isole alors \(P(A \cap B)\) :
\[
P(A \cap B) = P(B) - P(A^c \cap B) = 0{,}2 - 0{,}1 = 0{,}1.
\]

\[
\boxed{P(A \cap B) = 0{,}1}
\]    
\end{solution}


\question[10]

Douze personnes participent à une soirée micro-ouvert, dont Julie. Un programme de la soirée consiste à choisir \(5\) personnes parmi les \(12\) et à fixer un ordre de passage pour ces \(5\) personnes. Une contrainte est imposée :
\emph{la première personne à passer doit être l’une des trois personnes désignées comme animateurs}.
Julie n’est \emph{pas} animatrice.

Si on choisit un programme au hasard parmi tous les programmes respectant cette contrainte, quelle est la probabilité que Julie soit choisie et qu’elle passe en troisième position?

\begin{solution}

\textbf{Nombre total de programmes possibles}

On choisit d’abord l’animateur qui passe en premier : \(3\) choix.

Puis on choisit et ordonne \(4\) autres personnes parmi les \(11\) restantes :
\[
{}_{11}P_{4}.
\]

Le nombre total de programmes est donc
\[
3\times {}_{11}P_{4}
=3\times \frac{11!}{7!}.
\]

\medskip
\textbf{Cas favorables}

On impose que Julie soit choisie et qu’elle passe en troisième position.

\emph{Étape 1 :} choisir l’animateur en première position : \(3\) choix.

\emph{Étape 2 :} choisir les \(3\) autres personnes parmi les \(10\) restantes (en excluant Julie et l’animateur choisi) :
\[
\binom{10}{3}.
\]

\emph{Étape 3 :} ordonner ces \(3\) personnes et Julie dans les positions
\(2,3,4,5\), en imposant que Julie soit en position \(3\).

Il reste alors \(3!\) ordres possibles.

Le nombre de programmes favorables est donc
\[
3\times \binom{10}{3}\times 3!
\]

\medskip
\textbf{Probabilité}

\[
\mathbb{P}(\text{Julie choisie et en 3e})
=
\frac{3\,\binom{10}{3}\,3!}{3\,{}_{11}P_{4}}
=
\frac{\binom{10}{3}\,3!}{{}_{11}P_{4}}.
\]

On simplifie :
\[
\binom{10}{3}\,3!
=
\frac{10!}{7!},
\qquad
{}_{11}P_{4}=\frac{11!}{7!}.
\]

Donc
\[
\mathbb{P}(\text{Julie choisie et en 3e})
=
\frac{10!}{11!}
=
\frac{1}{11}.
\]

\[
\boxed{\mathbb{P}=\frac{1}{11}.}
\]
\end{solution}

%%%%%%%%%%%%%%%%%%% PROB COND. %%%%%%%%%%%%%%

\question[]

Un sac contient 100 dés à 6 faces, dont 60 sont rouges et 40 sont bleus. Les 60 dés rouges sont numérotés de <<1>> à <<6>> et sont réguliers (c'est-à-dire que les faces sont équiprobables). Les 40 dés bleus ont 6 fois la face avec le numéro <<1>>. L'expérience aléatoire consiste à choisir un dé dans le sac, à le lancer et à noter la face du dessus.

Pour répondre aux questions, utilisez $\mathcal{U}$ pour désigner l'espace échantillonnal et les événements suivants:

\begin{itemize}
    \item R: le dé choisi est rouge,
    \item B: le dé choisi est bleu,
    \item N: la face sur le dessus du dé lancé est <<1>>,
    \item S: la face sur le dessus du dé lancé est <<6>>.
\end{itemize}

\begin{parts}
    \part[5] La face sur le dessus du dé lancé est <<1>>. Quelle est la probabilité que le dé lancé soit rouge?


\begin{solution}
    On sait que $P(R)=0,6$, $P(B)=0,4$, $P(N|R)=1/6$, $P(N|B)=1$. On peut déduire que 

    \begin{align*}
        P(N)&=P(N\cap R)+P(N\cap B)\\
        &= P(N|R)P(R)+P(N|B)P(B)\\
        &=\frac{1}{6}\times 0,6 + 0,4 \times 1\\
        &=0,5
    \end{align*}

    On cherche $P(R|N)$.

    \begin{align*}
        P(R|N)&=\frac{P(N\cap R)}{P(N)}\\
        &=\frac{P(N|R)P(R)}{P(N)}\\
        &= \frac{\sfrac{1}{6}\times 0,6}{0,5}\\
        &=0,2
    \end{align*}
\end{solution}

\vfill

\part Dites si les énoncés suivants sont vrais ou faux.

\begin{subparts}
    \subpart[1] Les événements $R$ et $B$ sont indépendants.

\begin{oneparcheckboxes}
    \choice Vrai
    \correctchoice Faux
\end{oneparcheckboxes}
\begin{solution}
$P(R\cap B)=0$ car un dé ne peut être à la fois rouge et bleu, mais $P(R)P(B)=0{,}6\times 0{,}4=0{,}24\neq 0$.
\end{solution}
\vspace{0.3cm}

    \subpart[1] Les événements $R$ et $B$ sont mutuellement exclusifs.

\begin{oneparcheckboxes}
    \correctchoice Vrai
    \choice Faux
\end{oneparcheckboxes}
\begin{solution}
Un dé ne peut être à la fois rouge et bleu, donc $P(R\cap B)=0$.
\end{solution}
\vspace{0.3cm}

    \subpart[1] Les événements $B$ et $N$ sont indépendants.

\begin{oneparcheckboxes}
    \choice Vrai
    \correctchoice Faux
\end{oneparcheckboxes}
\begin{solution}
$P(B\cap N)=P(N|B)\,P(B)=1\times 0{,}4=0{,}4$, mais $P(B)\,P(N)=0{,}4\times 0{,}5=0{,}2\neq 0{,}4$.
\end{solution}
\vspace{0.3cm}

    \subpart[1] Les événements $B$ et $S$ sont mutuellement exclusifs.

\begin{oneparcheckboxes}
    \correctchoice Vrai
    \choice Faux
\end{oneparcheckboxes}
\begin{solution}
Un dé bleu n'a que des faces <<1>>, donc $P(S|B)=0$ et $P(B\cap S)=0$.
\end{solution}
\vspace{0.3cm}

    \subpart[1] Les événements $R$ et $S$ sont indépendants.

\begin{oneparcheckboxes}
    \choice Vrai
    \correctchoice Faux
\end{oneparcheckboxes}
\begin{solution}
$P(R\cap S)=P(S|R)\,P(R)=\frac{1}{6}\times 0{,}6=0{,}1$, mais $P(R)\,P(S)=0{,}6\times 0{,}1=0{,}06\neq 0{,}1$.
\end{solution}
\end{subparts}
\end{parts}

\newpage

\question[2]

On peut modéliser le nombre de jours avant la panne d'un serveur informatique par une loi géométrique de paramètre $p=0{,}4$. Après 10 jours sans panne, que peut-on dire?

\begin{checkboxes}
\choice La panne est maintenant plus probable qu'au début.
\choice La panne est moins probable qu'au début.
\correctchoice La probabilité de panne demain est la même qu'au premier jour.
\choice Il est certain qu'une panne surviendra bientôt.
\end{checkboxes}
\begin{solution}
La loi géométrique possède la propriété d'absence de mémoire~: $\mathbb{P}(X > s+t \mid X > s) = \mathbb{P}(X > t)$. Le fait d'avoir survécu 10 jours sans panne ne change pas la probabilité de panne le lendemain.
\end{solution}

\newpage 

%%%%%%%% Module 3 - lois discrètes %%%%%%%%%%%%%

\question[]

Dans un projet de cartographie, un drone tente d’envoyer des relevés de
position vers un serveur. À chaque tentative, le relevé est soit reçu
avec succès, soit perdu. On suppose que les tentatives sont indépendantes
et que la probabilité qu’un relevé soit reçu avec succès est constante et
égale à $p=0{,}6$.

Charles surveille la réception des relevés.

\begin{parts}

\part[4]
Au cours des $10$ prochaines tentatives, quelle est la probabilité que
\emph{exactement $7$} relevés soient reçus avec succès?

\begin{solution}
Soit $X$ le nombre de relevés reçus avec succès parmi $10$ tentatives.
Alors
\[
X \sim \text{Binomiale}(10,0{,}6).
\]
On a donc
\[
\mathbb P(X=7)
= \binom{10}{7}(0{,}6)^7(0{,}4)^3.
\]
\end{solution}

\part[4]
Quelle est la probabilité que le \emph{premier} relevé reçu avec succès
survienne à la $5^{\text{e}}$ tentative?

\begin{solution}
Soit $T$ le numéro de la tentative correspondant au premier succès.
Alors
\[
T \sim \text{Géométrique}(0{,}6).
\]
On obtient
\[
\mathbb P(T=5)=(1-0{,}6)^4(0{,}6)=(0{,}4)^4(0{,}6).
\]
\end{solution}

\part[4]
Chaque tentative d’envoi coûte $0{,}40\$$, qu’elle soit réussie ou non.
De plus, le lancement du drone entraîne un coût fixe de $10\$$.
On souhaite obtenir $50$ relevés reçus avec succès.

On note $Z$ le coût total (en dollars) nécessaire pour atteindre cet
objectif. Calculez $\mathbb E[Z]$.

\begin{solution}
Soit $Y$ le nombre de tentatives nécessaires pour obtenir $50$ relevés
reçus avec succès. Comme chaque tentative est indépendante et que la
probabilité de succès est $p=0{,}6$, on a
\[
Y \sim \text{Pascal}(50,0{,}6).
\]

Pour une loi de Pascal de paramètres $(r,p)$, l’espérance est
\[
\mathbb E[Y] = \frac{r}{p}.
\]
Ainsi,
\[
\mathbb E[Y] = \frac{50}{0{,}6} \approx 83{,}33.
\]

Le coût total s’écrit
\[
Z = 10 + 0{,}40\,Y.
\]
Par linéarité de l’espérance,
\[
\mathbb E[Z]
= 10 + 0{,}40\,\mathbb E[Y]
= 10 + 0{,}40 \times \frac{50}{0{,}6}.
\]
Numériquement,
\[
\mathbb E[Z] \approx 10 + 33{,}33 = 43{,}33\ \text{\$}.
\]
\end{solution}


\end{parts}

\newpage

%%%%%%% Module 4 partie a - lois continuesprocessus de poisson %%%%%%%%

\question[]

On s'intéresse au temps avant d'obtenir une contravention. On suppose que le temps $T_1$ (en jours) avant la première contravention
suit une loi exponentielle dont l’espérance est de $5$ jours. Les temps entre les contraventions sont indépendants et de même loi.

Chaque contravention coûte $65\$$ et un permis mensuel coûte $180\$$.

\begin{parts}

\part[8]
Considérons un mois de $30$ jours. Quelle est la probabilité que le coût total des contraventions dépasse le coût du permis mensuel?

\begin{solution}
Pour une loi exponentielle de paramètre $\lambda$, on a
\[
\mathbb E[T_1] = \frac{1}{\lambda}.
\]
Comme $\mathbb E[T_1]=5$, on obtient
\[
\lambda = \frac{1}{5} \text{ par jour}.
\]

Dans ce modèle, le nombre de contraventions sur un intervalle de longueur
$t$ suit une loi de Poisson de paramètre $\lambda t$.
Ainsi,
\[
N(30) \sim \text{Poisson}\!\left(\frac{1}{5}\times 30\right)
= \text{Poisson}(6).
\]

Le coût total dépasse $180\$$ dès qu’il y a au moins $3$ contraventions,
car $2\times 65 = 130 < 180$ et $3\times 65 = 195 > 180$.
Donc
\[
\mathbb P(N(30)\ge 3)
= 1 - \sum_{k=0}^{2} e^{-6}\frac{6^k}{k!}
= 1 - e^{-6}\left(1 + 6 + \frac{6^2}{2}\right).
\]
Numériquement, cette probabilité vaut environ $0{,}94$.
\end{solution}

\part[4]
Soit $T_3$ le temps (en jours) avant la troisième contravention. Calculez $\mathbb E[T_3]$.

\begin{solution}
Le temps $T_3$ est la somme de trois temps exponentiels indépendants
de paramètre $\lambda$, donc suit une loi gamma de forme $3$ et de taux
$\lambda$.

On a
\[
\mathbb E[T_3] = \frac{3}{\lambda}.
\]
Avec $\lambda = \frac{1}{5}$, on obtient
\[
\mathbb E[T_3] = 15 \text{ jours}.
\]
\end{solution}

\end{parts}


\newpage
%%%%%%%%%%%%%%%% Module 2 - Variable aléatoire continue %%%%%%%%%%%%%%%%

\question[] Soit $X$ une variable aléatoire continue dont la fonction de répartition est donnée par
\[ F_X(x) = \begin{cases}
\displaystyle 0 & \mbox{si } x < 0, \\[6pt]
\displaystyle cx^3 & \mbox{si } 0 \le x < 2, \\[6pt]
\displaystyle 1 & \mbox{si } x \ge 2.
\end{cases} \]

\begin{parts}
    \part[4] Trouvez la valeur de la constante $c$.

    \begin{solution}
    On utilise la condition $F_X(2)=1$:
    \[
    c(2)^3 = 1 \implies 8c = 1 \implies c = \frac{1}{8}.
    \]
    \end{solution}
    \vspace{4cm}

    \part[4] Trouvez la fonction de densité $f_X(x)$.

    \begin{solution}
    On dérive la fonction de répartition:
    \[
    f_X(x) = F_X'(x) = \begin{cases}
    \displaystyle \frac{3x^2}{8} & \mbox{si } 0 \le x < 2, \\[6pt]
    0 & \mbox{sinon.}
    \end{cases}
    \]
    \end{solution}
    \vspace{5cm}

    \part[4] Calculez $\mathbb{P}\!\left(X > 1 \;\middle|\; X < \dfrac{3}{2}\right)$.

    \begin{solution}
    \begin{align*}
    \mathbb{P}\!\left(X > 1 \mid X < \tfrac{3}{2}\right)
    &= \frac{\mathbb{P}\!\left(1 < X < \frac{3}{2}\right)}{\mathbb{P}\!\left(X < \frac{3}{2}\right)}
    = \frac{F_X\!\left(\frac{3}{2}\right) - F_X(1)}{F_X\!\left(\frac{3}{2}\right)}.
    \end{align*}

    On calcule:
    \[
    F_X\!\left(\tfrac{3}{2}\right) = \frac{(3/2)^3}{8} = \frac{27}{64}, \qquad
    F_X(1) = \frac{1}{8} = \frac{8}{64}.
    \]
    Donc
    \[
    \mathbb{P}\!\left(X > 1 \mid X < \tfrac{3}{2}\right)
    = \frac{\frac{27}{64} - \frac{8}{64}}{\frac{27}{64}}
    = \frac{19}{27}.
    \]
    \end{solution}
    \vspace{5cm}

    \part[4] Calculez $\mathbb{E}(X)$.

    \begin{solution}
    \begin{align*}
    \mathbb{E}(X)
    &= \int_0^2 x \cdot \frac{3x^2}{8}\,dx
    = \frac{3}{8}\int_0^2 x^3\,dx
    = \frac{3}{8}\left[\frac{x^4}{4}\right]_0^2
    = \frac{3}{8} \cdot \frac{16}{4}
    = \frac{3}{2}.
    \end{align*}
    \end{solution}
\end{parts}



\newpage
%%%%%%%%%%%%%%%% Module 4 - Loi normale et TCL %%%%%%%%%%%%%%%%

\question[] Un fabricant de céréales sait que la masse $X$ (en grammes) d'une boîte de céréales suit une loi normale de moyenne $\mu = 500$ g et de variance $\sigma^2 = 225$ g$^2$.

\begin{parts}
    \part[4] Quelle est la probabilité qu'une boîte choisie au hasard ait une masse entre $480$ g et $520$ g?

    \begin{solution}
    On a $X \sim \mathcal{N}(500, 225)$, donc $\sigma = 15$.
    \begin{align*}
    \mathbb{P}(480 < X < 520)
    &= \mathbb{P}\!\left(\frac{480-500}{15} < Z < \frac{520-500}{15}\right) \\
    &= \mathbb{P}(-1{,}33 < Z < 1{,}33) \\
    &= 2\,\Phi(1{,}33) - 1.
    \end{align*}
    \end{solution}
    \vspace{6cm}

    \part[6] On prélève un échantillon de $36$ boîtes. Quelle est la probabilité que la masse moyenne de ces $36$ boîtes dépasse $504$ g?

    \begin{solution}
    Par la loi de la moyenne échantillonnale,
    \[
    \overline{X} \sim \mathcal{N}\!\left(500,\;\frac{225}{36}\right) = \mathcal{N}(500;\; 6{,}25),
    \]
    donc $\sigma_{\overline{X}} = 2{,}5$.
    \begin{align*}
    \mathbb{P}(\overline{X} > 504)
    &= \mathbb{P}\!\left(Z > \frac{504 - 500}{2{,}5}\right)
    = \mathbb{P}(Z > 1{,}6)
    = 1 - \Phi(1{,}6).
    \end{align*}
    \end{solution}
    \vspace{8cm}

    \part[8] Le coût de production d'une boîte (en dollars) est une variable aléatoire dont on ne connaît pas la loi, mais dont l'espérance est $2$\$ et la variance est $0{,}64$\$$^2$. Approximez la probabilité que le coût total de production de $100$ boîtes dépasse $212$\$.

    \begin{solution}
    Soit $C_i$ le coût de la $i$-ème boîte. On a $\mathbb{E}(C_i)=2$ et $\text{Var}(C_i)=0{,}64$.

    Soit $T = C_1 + C_2 + \cdots + C_{100}$ le coût total.
    \[
    \mathbb{E}(T) = 100 \times 2 = 200, \qquad
    \text{Var}(T) = 100 \times 0{,}64 = 64.
    \]

    Par le théorème central limite ($n=100$ est assez grand),
    \[
    T \;\dot\sim\; \mathcal{N}(200,\;64), \quad \sigma_T = 8.
    \]
    Donc
    \begin{align*}
    \mathbb{P}(T > 212)
    &\approx \mathbb{P}\!\left(Z > \frac{212 - 200}{8}\right)
    = \mathbb{P}(Z > 1{,}5)
    = 1 - \Phi(1{,}5).
    \end{align*}
    \end{solution}
\end{parts}



\newpage
%%%%%%%%%%%%%%%% Module 5 - Loi conjointe %%%%%%%%%%%%%%%%

\question[] Soit $(X, Y)$ un couple de variables aléatoires discrètes dont la fonction de masse conjointe est donnée par le tableau suivant, où $a$ et $b$ sont des inconnues.

\begin{center}
\begin{tabular}{cc|cc}
$p(x,y)$ & & \multicolumn{2}{c}{$y$} \\
 & & $0$ & $1$  \\
\hline
 & $1$ & $\sfrac{1}{12}$ & $a$  \\
$x$ & $2$ & $b$ & $\sfrac{1}{6}$  \\
 & $3$ & $\sfrac{1}{4}$ & $\sfrac{1}{12}$  \\
\end{tabular}
\end{center}

On sait que $\mathbb{E}(X)=2$.

\begin{parts}
    \part[4] Que valent $a$ et $b$?

    \begin{solution}
    La somme des probabilités vaut $1$:
    \[
    \frac{1}{12} + a + b + \frac{1}{6} + \frac{1}{4} + \frac{1}{12} = 1
    \implies a + b = 1 - \frac{7}{12} = \frac{5}{12}.
    \]

    On utilise $\mathbb{E}(X)=2$. Les lois marginales de $X$ sont:
    \[
    \mathbb{P}(X=1)=\frac{1}{12}+a,\quad
    \mathbb{P}(X=2)=b+\frac{1}{6},\quad
    \mathbb{P}(X=3)=\frac{1}{4}+\frac{1}{12}=\frac{1}{3}.
    \]
    Donc
    \[
    \mathbb{E}(X) = 1\!\left(\frac{1}{12}+a\right) + 2\!\left(b+\frac{1}{6}\right) + 3\!\left(\frac{1}{3}\right) = \frac{17}{12} + a + 2b = 2.
    \]
    On obtient $a + 2b = \frac{7}{12}$. En soustrayant $a+b=\frac{5}{12}$, on trouve $b = \frac{2}{12} = \frac{1}{6}$, puis $a = \frac{5}{12} - \frac{1}{6} = \frac{1}{4}$.

    \[
    \boxed{a = \frac{1}{4}, \quad b = \frac{1}{6}.}
    \]
    \end{solution}
    \vspace{6cm}

    \part[6] Les variables $X$ et $Y$ sont-elles indépendantes? Justifiez votre réponse.

    \begin{solution}
    On vérifie si $\mathbb{P}(X=x, Y=y) = \mathbb{P}(X=x)\,\mathbb{P}(Y=y)$ pour tout $(x,y)$.

    Marginales:
    \[
    \mathbb{P}(X=1)=\frac{1}{12}+\frac{1}{4}=\frac{1}{3},\qquad
    \mathbb{P}(Y=0)=\frac{1}{12}+\frac{1}{6}+\frac{1}{4}=\frac{1}{2}.
    \]

    \[
    \mathbb{P}(X=1)\,\mathbb{P}(Y=0)=\frac{1}{3}\times\frac{1}{2}=\frac{1}{6},
    \]
    mais
    \[
    \mathbb{P}(X=1, Y=0)=\frac{1}{12}\neq \frac{1}{6}.
    \]

    Donc $X$ et $Y$ ne sont \textbf{pas indépendantes}.
    \end{solution}
    \vspace{6cm}

    \part[6] Calculez $\mathbb{E}[Y \mid X=2]$.

    \begin{solution}
    On a $\mathbb{P}(X=2) = b + \frac{1}{6} = \frac{1}{6}+\frac{1}{6} = \frac{1}{3}$.

    La loi conditionnelle de $Y$ sachant $X=2$ est:
    \[
    \mathbb{P}(Y=0 \mid X=2) = \frac{p(2,0)}{p_X(2)} = \frac{\sfrac{1}{6}}{\sfrac{1}{3}} = \frac{1}{2}, \qquad
    \mathbb{P}(Y=1 \mid X=2) = \frac{p(2,1)}{p_X(2)} = \frac{\sfrac{1}{6}}{\sfrac{1}{3}} = \frac{1}{2}.
    \]
    Donc
    \[
    \mathbb{E}[Y \mid X=2] = 0\times\frac{1}{2} + 1\times\frac{1}{2} = \frac{1}{2}.
    \]
    \end{solution}
\end{parts}



\end{questions}

\end{document}
